\section{Numerical replay and accepted support sensitivity}
The additional support analysis distinguishes numerical agreement with an archived prediction from stability of a scored metric. Using unchanged relative and absolute tolerances of $2\times10^{-5}$, 16 of 36 Alice neural replays pass; all nine TCN and nine recurrent cases plus Sanyo Inverted43/44 remain outside tolerance. The maximum difference in the saved mixed-origin replay is 1.287 W, with maximum whole-output difference RMS 0.02032 W. On identical scored targets at one and twelve hours, across full, active and low-power ranges, the largest RMSE change is 0.000567 W and the largest MAE change is 0.000425 W. No compared historical-support mean ranking or Daily-effect direction changes. These small metric effects do not prove that the historical neural input state has been recovered.

For common labels and mask, $|\mathrm{RMSE}(p)-\mathrm{RMSE}(q)|\leq\sqrt{\mathrm{mean}(p-q)^2}$ and $|\mathrm{MAE}(p)-\mathrm{MAE}(q)|\leq\mathrm{mean}|p-q|$. The recorded bounds pass on the checked supports. They do not constrain newly restored origins; relative skill can also be sensitive when the reference error is close to zero.

Representative probes cover Qcells TCN44, Sanyo TCN42, Sanyo Inverted43/44 and Qcells Inverted42. Strict checkpoint loading has no missing or unexpected parameter keys. Models are in evaluation mode with float32 inputs and frozen target inverse transforms. Original versus mixed 256-origin batches do not resolve the selected Alice discrepancies; single-item batches can change their size. Neither the CPU nor a different batch output is assumed to be the truth. The current runtime is PyTorch 2.7.1+cu118, CUDA 11.8, cuDNN 90100 and RTX 3060 Laptop GPU, with deterministic cuDNN, cuDNN TF32 enabled and matrix TF32 disabled. Complete historical backend settings and the original separately saved neural KNN/IF/scaler state are unavailable. The current reconstruction uses the later saved Train-only Ridge processor. Its feature order and fitted-state summaries are recorded, without claiming identity with an unavailable historical object.

\subsection{Per-method acceptance}
Hanwha and Qcells include all three verified Inverted seeds in new-support primary contrasts; Sanyo includes seed42 separately, never an incomplete three-seed mean. Joint-patch is verified secondary evidence at all three arrays. Original/expanded Ridge A/B and deterministic references replay successfully. Other new-origin neural metrics are clearly marked diagnostic and excluded from accepted averages/rankings. Full accepted and diagnostic CSVs are separate. This method-specific scope supersedes the earlier all-site scoring block while preserving its original records and tolerance.

For each prefix, original eligibility requires all targets in that prefix finite and nonnegative. S1 requires each scored target nonnegative; S2 retains each finite raw value. New origins are relative to the corresponding prefix's original eligibility, not always H144. Inputs and nonnegative Daily/Last-value rules remain unchanged, and all compared methods share the same Daily point intersection. A label beyond Test is unscored but does not invalidate an available previous-day input. No full-window labels are imputed or missing rows joined together. Counts include origins, origin--lead pairs, unique physical target timestamps and active composition. Train support remains unchanged: Qcells neural Train/Validation counts are 12747/2999, distinct from Ridge's 12648/2977.

\section{Four-cell period and support definitions}
Both external periods use every calendar origin and historical/target buffers: October--December 2017 and April--June 2018. January 1 records provide the final twelve-hour target buffer for December 31 origins; they do not add a new origin month. Complete-H144 and pointwise-finite masks share the same candidate list, finite origin requirement and Daily point intersection. NIST remains fixed EST without DST; Yulara retains its provider-local coordinate and conservative availability shift. Finite negative external values remain eligible.

All 2018 predictions are reused. The 2017 supplemented candidates use the frozen 2017 processors and checkpoints with no fits. Archival-origin forward passes reproduce the original external predictions. A Yulara mixed-candidate batch probe has a maximum 0.03624 W difference outside the old tolerance for one seed. The explicit inference partition therefore preserves original chronological batching for newly computed archival origins, then computes additional origins separately. Both sets contain new forwards from the same checkpoint; historical saved predictions are not spliced into a replacement predictor. The failed mixed-batch diagnostic is retained. Original archive counts and metrics are reported separately from unified-buffer comparisons.

\subsection{Paired uncertainty and independent checks}
Forty-eight-hour nonoverlapping origin blocks are the primary scheme, with 24/72-hour sensitivity, 2000 draws and seed 20260908. Methods and actual neural seeds use the same draw within each comparison. SSE and point counts are summed before RMSE; neural-seed effects are averaged, not predictions. Periods are analyzed separately rather than pairing different calendar dates. Empty blocks and zero-support draws remain explicitly counted. These intervals condition on frozen predictions and the observed periods and do not remove all dependence or include retraining uncertainty.

The current saved-point audit independently reconstructs masks, time joins and metrics, while a separate block-count-matrix implementation recomputes interval effects without importing the production interval function. Earlier numerical audit totals remain historical. New archive verification tests file integrity and lightweight arithmetic, not complete training reproduction.

